\chapter{Wstęp}\label{ch:wstep}

\TemplateNote{Rozdziały w katalogu \file{chapters/} są neutralnymi przykładami. Każdy z nich pokazuje inne elementy szablonu. Zastąp treść własną, a ramki z wskazówkami (makro \latexcmd{TemplateNote}) usuń.}

\section{Cel i zakres dokumentu}\label{sec:cel-zakres}

Celem niniejszego dokumentu jest przedstawienie przebiegu i wyników zrealizowanego zadania. W tym miejscu opisz
krótko motywację oraz problem, którym zajmuje się dokument. Zakres obejmuje w szczególności:
\begin{itemize}
    \item opis przyjętych założeń i wykorzystanych narzędzi,
    \item prezentację przebiegu realizacji zadania,
    \item analizę uzyskanych wyników wraz z wnioskami.
\end{itemize}

\section{Układ dokumentu}\label{sec:uklad}

Dokument składa się z sześciu rozdziałów. Rozdział~\ref{ch:przeglad} zawiera przegląd literatury i podstawy
teoretyczne, rozdział~\ref{ch:metodyka} opisuje przyjętą metodykę, a rozdział~\ref{ch:realizacja} -- przebieg
realizacji. Uzyskane wyniki omówiono w rozdziale~\ref{ch:wyniki}, natomiast całość podsumowuje rozdział~\ref{ch:podsumowanie}.
Odwołania do rozdziałów, sekcji, rysunków i tabel tworzy się poleceniem \latexcmd{ref\{etykieta\}}, zawsze
poprzedzonym twardą spacją (znak \texttt{\string~}), aby numer nie został oddzielony od słowa na końcu wiersza.

\section{Konwencje typograficzne}\label{sec:konwencje}

W całym dokumencie stosuje się jednolite makra do oznaczania elementów technicznych:
\begin{itemize}
    \item \latexcmd{code\{...\}} -- fragment kodu lub wyrażenie, np. \code{int licznik_zdarzen = 0;},
    \item \latexcmd{file\{...\}} -- nazwa pliku lub ścieżka, np. \file{/home/user/dane/wynik_koncowy.csv},
    \item \latexcmd{cmd\{...\}} -- polecenie konsoli, np. \cmd{sha256sum plik.bin},
    \item \latexcmd{hash\{...\}} -- długi ciąg znaków zawijany w dowolnym miejscu, np. skrót SHA-256:
          \hash{e3b0c44298fc1c149afbf4c8996fb92427ae41e4649b934ca495991b7852b855},
    \item \latexcmd{enquote\{...\}} -- cudzysłów w polskiej formie, np. \enquote{przykładowy cytat}.
\end{itemize}

Nagłówek wewnątrz sekcji (zamiast ręcznego pogrubiania) tworzy makro \latexcmd{lead\{...\}}:

\lead{Przykładowy nagłówek wewnątrz sekcji}
Treść następująca bezpośrednio po nagłówku. Dłuższe fragmenty dzielimy poleceniami
\latexcmd{section}, \latexcmd{subsection} i \latexcmd{subsubsection}.
